% Ad Atticum 10.11 --- headnote
% ad-atticum-10-11, 4 May 49 BC, Cumae

\textit{Cicero to Atticus, written from the Cuman villa
on the fourth day before the Nones of May 49 BC --- 4
May (the manuscript dateline: \textit{Scr.\ in Cumano iv
Non.\ Mai.\ a.\ 705 (49)}). The previous letter (10.10)
had been sealed but held back; Philotimus then arrived
with a letter from Atticus, and the present one is the
reply. The bulk of the letter is family business between
the two brothers and their dependants, set against the
financial squeeze of the civil war. On Quintus: Cicero
defends him --- nothing in his brother is festering or
two-faced (\textit{[Greek: hypoulon]}), nothing that one
conversation could not put right; he loves all his
people, Cicero himself most. On Quintus's son: indulged,
yes, and turned a bit fierce and arrogant by it, but the
deeper faults are his own root and not from Cicero's
softness. The money paragraph is a quick, frank window
onto the cash crunch: Q.~Axius will not pay back the
thirteen thousand sesterces Cicero lent his son,
pleading the times; Lepta and others do the same;
Cicero hears Quintus complain of being pressed for a
twenty-thousand debt and wonders at it --- but defends
him as neither slow nor tight-fisted by nature.}

\textit{Section 4 returns to the immediate crisis.
Antony arrived at the neighbouring villa yesterday
evening; an interview is imminent, though Antony may
just send word, and Cicero is now doing everything
covertly. The escape plan --- a small boat from the Bay
of Naples to Sicily or Malta --- pulls at him: he
remembers being sick with anxiety even on an undecked
Rhodian galley (\textit{[Greek: aphrakt\=o(i)]}) in
summer; what will a tiny skiff in early May do to him?
Trebatius was there and reported ``monsters'' from
Caesar's camp, including the possibility of Balbus
(Caesar's Spanish banker-agent, not a senator) entering
the Senate. The closing paragraph is an apologetic
postscript about a stiff exchange of letters with the
banker Vettienus: Vettienus had written peremptorily
(\textit{[Greek: apotom\=os]}) about Cicero's money,
Cicero wrote back with too much heat
(\textit{[Greek: thumik\=oteron]}) and addressed him as
``master of the mint'' because Vettienus had styled him
``consul''; Atticus is asked to smooth it over.}
